% Part:sets-functions-relations
% Chapter: sets
% Section: schroder-bernstein
% Hindi translation (hi-Deva-IN) of Open Logic commit 9620cc73f9c8e0ad003c514a5d3748f29611c4c0.

\documentclass[../../../include/open-logic-section]{subfiles}

\begin{document}

\olfileid[hi]{sfr}{siz}{sb}

\olsection{आकार की धारणा और श्रोडर--बर्नस्टीन}

\begin{explain}
यह एक सहज विचार है: यदि $A$, $B$ से बड़ा नहीं है और $B$, $A$ से बड़ा नहीं
है, तो $A$ और $B$ समसंख्यक हैं। सच कहें तो यदि यह विचार \emph{गलत} होता,
तो इस धारणा का औचित्य सिद्ध करना लगभग असंभव होता कि समसंख्यता की हमारी
परिभाषित धारणा का समुच्चयों के ``आकारों'' की तुलना से कोई संबंध है!{}
सौभाग्य से यह सहज विचार सही है। श्रोडर--बर्नस्टीन प्रमेय इसका औचित्य सिद्ध
करता है।
\end{explain}

\begin{thm}[श्रोडर--बर्नस्टीन]
	\ollabel{thm:schroder-bernstein}
	यदि $\cardle{A}{B}$ और $\cardle{B}{A}$,
	तो $\cardeq{A}{B}$।
\end{thm}

\begin{explain}
दूसरे शब्दों में, यदि $A$ से~$B$ में कोई !!a{injection} और $B$ से~$A$ में
कोई !!a{injection} है, तो $A$ से~$B$ में कोई !!a{bijection} है।

लेकिन इस परिणाम को सिद्ध करना वास्तव में बहुत \emph{कठिन} है। वास्तव में
कैंटर ने परिणाम कहा था, पर इसे दूसरों ने सिद्ध किया।\footnote{इतिहास के
विषय में अधिक जानकारी के लिए, उदाहरणार्थ,
\citet[pp.~165--6]{Potter2004} देखें।}
\oliflabeldef{sfr:cardinals:card-sb:sec}{हम श्रोडर--बर्नस्टीन को
\olref[sfr][cardinals][card-sb]{sec} में जाकर ही \emph{सिद्ध} कर पाएँगे।}{}%
अभी के लिए आपको इस पर विश्वास करना पड़ता है।

सौभाग्य से श्रोडर--बर्नस्टीन \emph{सही} है, और इससे यह पुष्ट होता है कि
हमारे परिभाषित संबंधों, अर्थात $\cardeq{A}{B}$ और $\cardle{A}{B}$, का
``आकार'' से संबंध है। इसके अतिरिक्त श्रोडर--बर्नस्टीन बहुत \emph{उपयोगी}
है। दो समसंख्यक समुच्चयों के बीच किसी !!a{bijection} की कल्पना करना कठिन हो
सकता है। श्रोडर--बर्नस्टीन प्रमेय तुलना को दो भागों में बाँट देता है, ताकि
हमें केवल पहले समुच्चय से दूसरे में कोई !!a{injection} और फिर उलटी दिशा में
कोई ऐसा फलन सोचना पड़े।
\end{explain}
\end{document}
